\title{Direct Preference Optimization: Your Language Model is Secretly a Reward Model}

\begin{abstract}
Large-scale unsupervised language models (LMs) acquire extensive knowledge about the world and develop certain reasoning capabilities. However, their fully unsupervised training makes it challenging to exert fine-grained control over their behavior. Current strategies aimed at enhancing the steerability of such models involve gathering human feedback in the form of relative judgments on the quality of generated outputs, subsequently refining the unsupervised LM based on these preferences—commonly employing reinforcement learning from human feedback (RLHF). Yet, RLHF entails a multi-step, often fragile process: first constructing a reward model aligned with human judgments, then employing reinforcement learning to adjust the LM in a way that increases the reward, while ensuring the LM does not diverge excessively from its original form. This work presents a novel reparameterization of the reward model within the RLHF pipeline, which facilitates an exact, closed-form derivation of the corresponding optimal policy. This innovation allows us to address the RLHF optimization objective using nothing more than a straightforward classification loss. Our proposed method, termed \textit{Direct Preference Optimization} (DPO), offers a stable and efficient alternative, dispensing with the need for language model sampling or extensive hyperparameter search during fine-tuning. Experimental results demonstrate that DPO enables LMs to conform to human preferences as effectively as, or surpassing, previous approaches. Importantly, DPO-based fine-tuning not only outperforms PPO-driven RLHF in sentiment control tasks but also equals or betters the performance in tasks like summarization and single-turn dialogue generation, all while being markedly simpler to implement and train.
\end{abstract}

\section{Introduction}
Large-scale unsupervised language models (LMs) trained on massive corpora exhibit a range of unexpected abilities~\citep{chowdhery2022palm, brown2020language, touvron2023llama, bubeck2023sparks}. Despite these advances, the data fueling their development is produced by humans whose objectives, interests, and competencies vary widely. Not all of these characteristics are desirable for LMs to replicate. For example, while it is beneficial for an AI coding assistant to \textit{recognize} common programming errors to facilitate their correction, we would prefer that, during code generation, the model emphasizes the (potentially infrequent) examples of expert coding practices found in its training data. In another instance, it is useful for an LM to be \textit{informed} about a widespread misconception shared by half the population; nevertheless, we do not wish for the model to affirm this misconception in 50\% of relevant outputs. Consequently, it is essential to distinguish and select the model’s \emph{target behaviors and outputs} from its extensive \textit{knowledge base and competencies} to create AI systems that are reliable, effective, and responsive to control~\citep{ouyang2022training}. Although reinforcement learning (RL) is commonly employed to align LMs with human values, we demonstrate that the RL-based learning objective can, in fact, be realized exactly using a straightforward binary cross-entropy loss, thereby streamlining the overall preference learning process.

From a broad perspective, current approaches shape language model behaviors by leveraging carefully assembled datasets reflecting human evaluative judgments regarding safety and usefulness. This stage of preference learning follows the initial unsupervised pre-training on large-scale textual data. While supervised fine-tuning based on human-authored examples of high-quality responses represents the most direct methodology, reinforcement learning from human or artificial feedback (RLHF/RLAIF; \citep{christiano2017deep,bai2022constitutional}) has proven to be the most impactful class of methods. RLHF involves fitting a reward model to a corpus of human preference data and subsequently optimizing the policy of the language model via RL to encourage responses with high reward, all while avoiding significant deviations from the original model parameters. Despite their capacity to produce highly capable conversational and coding systems, RLHF pipelines entail greater complexity compared to standard supervised learning, requiring the training of multiple models and repeated sampling from the model policy during training—thus introducing considerable computational overhead.

This work demonstrates a method for optimizing language models to align with human preferences directly, circumventing the need for explicit reward modeling or reinforcement learning techniques. We introduce \textit{{Direct Preference Optimization} (DPO)}, an approach that implicitly targets the same objective as standard RLHF methods—specifically, maximizing reward subject to a KL-divergence constraint—while maintaining ease of implementation and training. DPO operates by increasing the log-likelihood ratio in favor of preferred responses relative to less favored ones, employing adaptive, example-specific weighting to mitigate the model collapse frequently observed when using a straightforward probability ratio approach. As with prior methods, DPO assumes a theoretical preference model, such as the Bradley-Terry framework (\cite{bradley1952rankanalysis}), to evaluate the correspondence between a reward function and observed human preference data. In contrast to traditional pipelines that use this preference model to train a reward function before fitting a policy, DPO employs a variable transformation to directly express the preference loss in terms of the policy itself. Provided a collection of human-annotated preference pairs, DPO leverages a binary cross entropy loss to train the policy directly, yielding the policy that is optimal under a latent reward function derived from preference observations.

The principal contribution of this paper is the Direct Preference Optimization (DPO) algorithm, a straightforward, reinforcement-learning-free technique for preference-based language model training. Empirical results indicate that DPO matches or surpasses the performance of existing algorithms—including those based on PPO-driven RLHF—across a range of applications such as sentiment control, summarization, and conversation modeling, even with models as large as 6B parameters.

Progressively larger self-supervised language models demonstrate the capacity to perform certain tasks without labeled data through zero-shot learning \citep{radford2019language}, or by leveraging a limited number of examples via few-shot prompting \citep{gpt3,megatron,chowdhery2022palm}. Nevertheless, their capabilities on downstream tasks and adherence to user instructions are often substantially enhanced by fine-tuning these models using datasets comprised of instructions paired with human-generated responses \citep{mishra-etal-2022-cross,sanh2022multitask,chung2022scaling,thoppilan2022lamda}. This process, known as `instruction-tuning', not only improves the model's utility but also allows it to extend its abilities beyond the instructions explicitly present in the fine-tuning data \citep{chung2022scaling}. In contrast to collecting expert demonstrations, gathering \textit{relative} human evaluations of response quality tends to be more practical, prompting more recent efforts to fine-tune LLMs using datasets reflecting human preferences. This has resulted in marked gains in fields such as translation \citep{kreutzer-etal-2018-reliability}, summarization \citep{stiennon2022learning,ziegler2020finetuning}, story-telling \citep{ziegler2020finetuning}, and following instructions \citep{ouyang2022training,ramamurthy2023is}. The prevailing methodology involves initially training a neural network to approximate a reward function compatible with observed preference data, utilizing a preference model like the Bradley-Terry model \citep{bradley1952rankanalysis}. Subsequently, the language model is fine-tuned to maximize this learned reward by employing reinforcement learning techniques such as REINFORCE \citep{williams1992reinforce}, proximal policy optimization (PPO; \cite{schulman2017proximal}), or related algorithms \citep{ramamurthy2023is}. Closely connected research utilizes instruction-following LLMs, fine-tuned through human feedback, to produce further synthetic preference data focused on specific dimensions like safety or harmlessness \citep{bai2022constitutional}, relying on human-authored rubrics as a form of weak supervision for model-generated evaluations. These strategies merge advances from two domains: one that seeks to train language models through reinforcement learning for multiple objectives~\citep{Ranzato2015SequenceLT,paulus2018a,wu2018learning}, and another centered on developing general approaches to learn from human preferences \citep{christiano2017deep,kupcsik2018learning}. While harnessing relative human preferences is advantageous, the process of fine-tuning large language models via reinforcement learning presents significant technical barriers; this work introduces a theoretically grounded framework for directly optimizing with relative preferences, bypassing the need for RL.

Beyond language applications, policy learning from preferences has been explored within both bandit and reinforcement learning frameworks, resulting in a variety of methodologies. In the contextual bandit setting, approaches that rely on action preferences or rankings rather than explicit rewards are referred to as contextual dueling bandits (CDB; \cite{yue2012karmed,dudik2015contextual}). When absolute reward signals are unavailable, the theoretical study of CDBs replaces the optimal policy with the concept of a \textit{von Neumann winner}: a policy that is expected to outperform any competing policy at least 50\% of the time \citep{dudik2015contextual}. Notably, CDB approaches generally assume preference labels are provided in an online fashion. In contrast, preference-based learning from human annotators often relies on a fixed, offline dataset of preference-labeled action pairs \citep{yan2022human}. In a related paradigm, \textit{preference-based RL} (PbRL), learning is driven by binary preference signals from an unspecified 'scoring' function instead of explicit rewards \citep{BusaFekete2014,ruiz2023dueling}. Various PbRL algorithms have been devised, some of which support off-policy reuse of collected preference data. However, most such algorithms entail a two-stage process: they first learn an explicit model of the latent reward or scoring function and then optimize the policy with respect to this model \citep{jain2013learning,BusaFekete2014,christiano2017deep,sadigh2017active,kupcsik2018learning}. By contrast, our approach formulates a single-stage policy optimization procedure that directly seeks to satisfy given preferences.

\section{Preliminaries}\label{section:prelims}

Here, we provide an overview of the RLHF methodology as described in \citeauthor{ziegler2020finetuning} and further explored in subsequent work \citep{stiennon2022learning, bai2022training, ouyang2022training}. This pipeline conventionally unfolds over three main stages: (1) supervised fine-tuning (SFT); (2) preference collection and reward model training; and (3) reinforcement learning-based policy optimization.

\textbf{SFT}: The RLHF workflow typically initiates with supervised adaptation of a pre-trained language model. This is performed using curated, high-quality data corresponding to the specific downstream application (e.g., dialogue, summarization), resulting in the supervised fine-tuned policy $\pi^\text{SFT}$.

\textbf{Reward Modelling Phase}: In the subsequent stage, the SFT-tuned model is prompted with input $x$ to produce two completions, $(y_1, y_2)\sim \pi^\text{SFT}(y \mid x)$. Human annotators are then presented with these outputs and asked to indicate their preference between them, denoted $y_w\succ y_l \mid x$—where $y_w$ is the favored answer, and $y_l$ is the less preferred, given the prompt $x$. The underlying assumption is that such preferences are governed by some unobserved reward function $r^*(y, x)$, which is not directly accessible. Several probabilistic models for preferences have been considered, with the Bradley-Terry (BT) model \cite{bradley1952rankanalysis} being a prevalent example (although broader frameworks such as the Plackett-Luce model \citep{plackett1975analysis, luce2012individual} can also be incorporated if multiple answers are ranked). The BT model represents the human preference distribution $p^*$ in the following form:

\begin{equation}\label{eq:bradley-terry}
    p^*(y_1\succ y_2 \mid x)=\frac{\exp\left(r^*(x, y_1)\right)}{\exp\left(r^*(x, y_1)\right) + \exp\left(r^*(x, y_2)\right)}.
\end{equation}

Given access to a fixed set of preference data $\mathcal{D}=\bigl\{x^{(i)}, y_w^{(i)}, y_l^{(i)}\bigr\}_{i=1}^N$ drawn from $p^*$, one can represent the reward using a parametric model $r_{\phi}(x, y)$ and fit its parameters using the maximum likelihood method. Casting this task as binary classification, the training objective takes the form of a negative log-likelihood loss:

\begin{equation}\label{eq:reward_model}
    \mathcal{L}_R(r_{\phi}, \mathcal{D}) = -\mathbb{E}_{(x, y_w, y_l)\sim \mathcal{D}}\bigl[\log \sigma(r_{\phi}(x, y_w)- r_{\phi}(x, y_l))\bigr]
\end{equation}

where $\sigma$ denotes the logistic sigmoid. In language model applications, the reward network $r_{\phi}(x, y)$ is commonly initialized with the weights from a supervised fine-tuning (SFT) model $\pi^\text{SFT}(y \mid x)$, with a single linear head appended to the final transformer representation, generating a scalar reward prediction \cite{ziegler2020finetuning}. To limit reward variance, earlier studies typically normalize the reward outputs so that $\mathbb{E}_{x,y\sim \mathcal{D}}\left[r_\phi(x, y)\right] = 0$ for all prompts $x$.

\textbf{RL Fine-Tuning Phase}: In the reinforcement learning stage, the acquired reward function is employed to supply feedback to the language model. Building on previous research~\citep{jaques2017sequence, jaques2020human}, the optimization problem is expressed as

\begin{equation}\label{eq:RL}
\max_{\pi_{\theta}}  \mathbb{E}_{x\sim \mathcal{D}, y\sim \pi_{\theta}(y \mid x)}\bigl[r_{\phi}(x, y)\bigr] - \beta\mathbb{D}_{\textrm{KL}}\bigl[\pi_{\theta}(y\mid x)\mid \mid \pi_\text{ref}(y\mid x)\bigr],
\end{equation}

with $\beta$ serving to regularize departures from the reference policy $\pi_\text{ref}$, which is the original SFT model $\pi^\text{SFT}$. 
The language model policy $\pi_\theta$ is usually initialized with parameters from $\pi^\text{SFT}$ as well. This KL-regularization is vital; it ensures the model remains close to distributions where the reward model's predictions remain valid and safeguards diversity in generated outputs, avoiding degeneration to deterministic, singular responses. Because language is inherently discrete, this objective is not amenable to gradient-based methods and must be tackled using reinforcement learning techniques. The prevailing methodology \citep{ziegler2020finetuning, stiennon2022learning, bai2022training, ouyang2022training} reformulates the reward as ${r(x, y) = r_{\phi}(x, y) -\beta (\log \pi_{\theta}(y\mid x) - \log \pi_\text{ref}(y\mid x))}$, then applies Proximal Policy Optimization (PPO) \cite{schulman2017proximal} to maximize this objective.

\section{Direct Preference Optimization}\label{sec:DPO}

Given the practical difficulties in using standard reinforcement learning on large-scale tasks such as language model adaptation, our objective is to devise a straightforward, direct method for policy optimization from preferences. In contrast to traditional RLHF pipelines that involve learning a reward function and subsequently optimizing it through RL, our strategy exploits a specific way of parameterizing the reward function that makes it possible to identify the associated optimal policy analytically, bypassing the need for an RL loop. As will be shown, the crux of our method is an explicit, analytic correspondence between reward functions and optimal policies; this insight allows the reformulation of loss functions defined over rewards into equivalent losses directly over policies. This variable transformation sidesteps the need for a separate reward model, yet remains compatible with prevalent preference modeling frameworks like the

\textbf{DPO Objective Derivation.} We begin from the same reinforcement learning framework as established in previous works (see Eq.~\ref{eq:RL}), using a generic reward function $r$. In line with the literature~\citep{peters2007reinforcement, peng2019advantage, korbak2022reinforcement, go2023aligning}, it can be readily demonstrated that the optimal policy for the KL-regularized reward maximization problem described in Eq.~\ref{eq:RL} is given by:

\begin{equation}\label{eq:op_policy}
    \pi_r(y\mid x) = \frac{1}{Z(x)}\pi_\text{ref}(y\mid x)\exp\left(\frac{1}{\beta}r(x, y)\right),
\end{equation}

where the normalization constant $Z(x)$ is defined as $Z(x) =\sum_{y}\pi_\text{ref}(y\mid x)\exp\left(\frac{1}{\beta}r(x, y)\right)$. For detailed steps, refer to Appendix \ref{app:derivation1}. Notably, even when $r_\phi$ is the maximum likelihood estimate for the true reward $r^*$, directly estimating $Z(x)$ remains computationally challenging~\citep{korbak2022reinforcement, go2023aligning}, thus making this form difficult to exploit in practical settings. Nevertheless, Eq.~\ref{eq:op_policy} can be algebraically rearranged to rewrite the reward in terms of the optimal policy $\pi_r$, baseline policy $\pi_\text{ref}$, and the partition function $Z(\cdot)$. Taking the logarithm of both sides and rearranging yields:

\begin{equation}\label{eq:main_eq}
    r(x,y) =\beta \log \frac{\pi_r(y\mid x)}{\pi_\text{ref}(y\mid x)} + \beta \log Z(x).
\end{equation}

This reparameterization can similarly be used for the true reward $r^*$ and its associated optimal policy $\pi^*$. The Bradley-Terry framework evaluates preference based solely on reward differences for two options: ${p^*(y_1 \succ y_2 \mid x) = \sigma(r^*(x, y_1) - r^*(x, y_2))}$. If we substitute Eq.~\ref{eq:main_eq} (for $r^*(x,y)$) into the preference probability defined in Eq.~\ref{eq:bradley-terry}, the term involving the partition function cancels. This allows us to rewrite the human preference probability directly in terms of $\pi^*$ and $\pi_\text{ref}$, without explicit dependence on the reward function. Therefore, under the Bradley-Terry model, the optimal RLHF policy $\pi^*$ fulfills:

\begin{equation}\label{eq:objective}
    p^*(y_1\succ y_2 \mid x)=\frac{1}{1 + \exp\left(\beta \log \frac{\pi^*(y_2\mid x)}{\pi_\text{ref}(y_2\mid x)} - \beta \log \frac{\pi^*(y_1\mid x)}{\pi_\text{ref}(y_1\mid x)}\right)}
\end{equation}

Details of the derivation are in Appendix~\ref{app:derivation2}. Although the above relies on the Bradley-Terry model, an analogous treatment applies for more general ranking models, such as the Plackett-Luce family~\citep{plackett1975analysis, luce2012individual}, discussed in Appendix~\ref{app:plackett_luce_models}.

Having cast the human preference likelihood in terms of the policy itself, we are now positioned to establish a maximum likelihood training criterion for a parameterized policy $\pi_\theta$. In parallel to the strategy employed in reward model training (cf. Eq.~\ref{eq:reward_model}), we arrive at the following policy objective:

\begin{equation}\label{eq:optimum_model}
    \mathcal{L}_\text{DPO}(\pi_{\theta}; \pi_\text{ref}) = -\mathbb{E}_{(x, y_w, y_l)\sim \mathcal{D}}\left[\log \sigma \left(\beta \log \frac{\pi_{\theta}(y_w\mid x)}{\pi_\text{ref}(y_w\mid x)} - \beta \log \frac{\pi_{\theta}(y_l\mid x)}{\pi_\

In this framework, we learn an implicit reward by adopting a different parameterization, with the optimal policy directly represented by $\pi_\theta$. Notably, the method corresponds to fitting a reparameterized Bradley-Terry model, which guarantees several theoretical properties—such as consistency—under appropriate assumptions about how the preference data are distributed \cite{bong2022generalized}. We examine additional theoretical aspects of DPO and how it compares to related approaches in Section~\ref{sec:theory}.

\textbf{How does the DPO update operate?} To gain a mechanistic insight into DPO, one can investigate the gradient of the objective function $\mathcal{L}_\text{DPO}$. The gradient with respect to the parameters $\theta$ takes the following form:

\begin{multline*}\label{eq:gradient}
    \nabla_\theta \mathcal{L}_\text{DPO}(\pi_\theta;\pi_\text{ref}) = \\ -\beta\mathbb{E}_{(x, y_w, y_l) \sim \mathcal{D}} \bigg[\underbrace{\sigma(\hat{r}_\theta(x, y_l) - \hat{r}_\theta (x, y_w))}_\text{greater emphasis when reward is misestimated}\bigg[\underbrace{\nabla_\theta\log \pi(y_w \mid x)}_\text{boost probability of $y_w$} - \underbrace{\nabla_\theta\log\pi(y_l \mid x)}_\text{diminish probability of $y_l$}\bigg]\bigg],
\end{multline*}

Here, $\hat{r}_\theta(x, y) = \beta \log \frac{\pi_\theta(y \mid x)}{\pi_\text{ref}(y \mid x)}$ defines the reward assigned by the language model $\pi_\theta$ relative to the reference $\pi_\text{ref}$ (with further discussion in Section~\ref{sec:theory}). Conceptually, the gradient of $\mathcal{L}_\text{DPO}$ works by raising the probability of the favored outputs $y_w$ and reducing the probability of the less preferred outputs $y_l$. Crucially, each training example is weighted by the degree to which the implicit reward model $\hat{r}_\theta$ favors the dispreferred over the preferred completion—modulated by $\beta$—thus reflecting the extent of misordering by the reward model as well as the influence of the KL regularization. Our empirical results indicate that this weighting plays a vital role; omitting the weighting term and using a simpler approach can lead to degeneration of the language model (see Appendix Table~\ref{tab:unlikelihood_generations}).

\textbf{DPO pipeline overview.} The standard workflow for DPO consists of the following steps: 1) For each prompt $x$, generate completions $y_1, y_2$ from $\pi_\text{ref}(\cdot \mid x)$, and then assign preference labels based on human feedback to create the offline preference dataset $\mathcal{D} = \{x^{(i)}, y_w^{(i)}, y_l^{(i)}\}_{i=1}^N$; 2) Update the language model $\pi_\theta$ by minimizing $\mathcal{L}_\text{DPO}$, conditioned on the chosen $\pi_\text{ref}$, the dataset $\mathcal{D}$, and the target value of $\beta$. In practice, it is desirable to leverage publicly available preference datasets instead of generating new samples and collecting human annotations. Since such datasets are typically drawn from $\pi^\text{SFT}$, we set $\pi_\text{ref} = \pi^\text{SFT}$ when possible. If $\pi^\text{SFT}$ is not provided, we obtain $\pi_\text{ref}$ by fitting it via maximum likelihood estimation to preferred responses ${(x, y_w)}$ in the data, i.e., ${\pi_\text{ref} = \argmax_{\pi}\mathbb{E}_{x, y_w \sim \mathcal{D}}\left[\log \pi(y_w \mid x)\right]}$. This step helps reduce distributional discrepancies between the true but inaccessible reference distribution and the practical $\pi_\text{ref}$ employed within DPO. Additional implementation details and hyperparameter settings are included in Appendix~\ref{app:implementation}.

\section{Theoretical Analysis of DPO} This section offers a deeper theoretical understanding of DPO, elucidates its foundation, and highlights its strengths relative to traditional actor-critic RLHF approaches such as PPO~\cite{schulman2017proximal}.

\label{sec:theory}

\subsection{Your Language Model Is Secretly a Reward Model} DPO enables the simultaneous learning of the policy without the need for explicit reward modeling or a separate RL optimization, instead utilizing a single maximum likelihood framework. The objective in Eq. \ref{eq:main_eq} can be interpreted as a Bradley-Terry model with a particular reward structure given by $r^*(x, y) = \beta \log\frac{\pi^*_\theta(y \mid x)}{\pi_\text{ref}(y \mid x)}$. The parametric model $\pi_\theta$ is thus trained analogously to optimizing a reward model, as formalized in Eq. \ref{eq:reward_model}, following an appropriate change of variables. The remainder of this section constructs the underlying theory supporting this reparameterization, demonstrating that it does not restrict the set of representable reward models and enables precise recovery of the optimal policy. We begin by introducing an equivalence relation over reward functions.

\begin{definition}
Two reward functions $r(x, y)$ and $r'(x, y)$ are considered equivalent if and only if there exists some function $f$ such that ${r(x, y)-r'(x, y) = f(x)}$.
\end{definition}

It is straightforward to verify that this defines an equivalence relation, which divides the collection of reward functions into distinct equivalence classes. This observation leads to the following two lemmas:

\begin{lemma}\label{lemma:same_prefrence}
Within the framework of the Plackett-Luce model, and specifically under the Bradley-Terry model, any two reward functions belonging to the same equivalence class generate identical preference distributions.
\end{lemma}

\begin{lemma}\label{lemma:same_policy}
Any pair of reward functions from a shared equivalence class produce the same optimal policy when solving the corresponding constrained RL problem.
\end{lemma}

Since the proofs of these statements are elementary, we provide them in Appendix \ref{app:lemma1}. The first lemma illustrates a classical identifiability ambiguity in Plackett-Luce-type models \cite{plackett1975analysis}. Because of this inherent ambiguity, it is typical to impose supplementary constraints to guarantee identifiability of the maximum likelihood estimator as in Eq. \ref{eq:reward_model} \cite{bong2022generalized}. The second lemma demonstrates that all reward functions in a given class yield equivalent optimal policies; thus, our learning goal reduces to finding any representative of the optimal class. We formally prove the following theorem in Appendix~\ref{app:thm1}:

\begin{theorem}\label{thm:main}
    Provided certain mild conditions, every reward class compatible with the Plackett-Luce (or specifically the Bradley-Terry) models admits the representation ${r(x, y) = \beta \log \frac{\pi(y\mid x)}{\pi_\text{ref}(y\mid x)}}$ for some policy model $\pi(y\mid x)$ and a fixed reference policy $\pi_\text{ref}(y \mid x)$.
\end{theorem}

\begin{sproof}
    Take any reward function $r(x, y)$, which specifies an associated optimal model $\pi_r(y \mid x)$ via Eq. \ref{eq:op_policy}. We aim to show that it is always possible to express a reward function from the equivalence class of $r$ using the stated reparameterization. Define the mapping $f$ as follows:
\begin{equation}
    f(r; \pi_\text{ref}, \beta)(x, y) = r(x, y) - \beta\log\sum_{y}\pi_\text{ref}(y\mid x)\exp\left(\frac{1}{\beta}r(x, y)\right)
\end{equation}
This operator $f$ normalizes the reward function by subtracting the logarithm of the partition function associated with $\pi_r$. Since the normalization depends solely on the prefix $x$, the function $f(r; \pi_\text{ref}, \beta)(x, y)$ is a member of the equivalence class of $r(x, y)$. Now, substituting $r$ using the right-hand side of Eq.~\ref{eq:main_eq} (which is valid for any reward), we find $f(r; \pi_\text{ref}, \beta)(x, y) = \beta \log \frac{\pi_r(y\mid x)}{\pi_\text{ref}(y\mid x)}$. Consequently, the mapping $f$ yields a function with the required parametric form from within the equivalence class, demonstrating that adopting this reparameterization does not diminish the generality of our reward modeling framework.
\end{sproof}

An alternative interpretation of Theorem~\ref{thm:main} is that it explicitly determines the particular reward function chosen within each equivalence class by the DPO reparameterization—specifically, the reward function that fulfills the following condition:

\begin{equation}\label{eq:lag_p}
     \sum_{y}\underbrace{\pi_\text{ref}(y\mid x)\exp\left(\frac{1}{\beta}r(x, y)\right)}_{=\pi(y\mid x)\text{, using Thm.~\ref{thm:main} reparam.}} = 1,
\end{equation}

This ensures that $\pi(y\mid x)$ forms a proper probability distribution (all probabilities are non-negative and their sum equals 1). Moreover, referencing Eq.~\ref{eq:op_policy}, Eq.~\ref{eq:lag_p} can be interpreted as the partition function associated with the optimal policy generated by the reward $r(x, y)$. The core idea of the DPO method is to introduce appropriate constraints on the generally under-specified Plackett-Luce (and specifically, the Bradley-Terry) family of preference models. This preserves the set of reward functions that can be represented, while ensuring that the optimal policy, as stated in Eq.~\ref{eq:op_policy}, remains tractable in closed form for any prompt $x$.

\subsection{Instability of Actor-Critic Algorithms} Our framework also provides insight into the sources of instability in common actor-critic approaches for RLHF, such as PPO. In keeping with the RLHF methodology, we focus on the RL fine-tuning procedure discussed in Section \ref{section:prelims}. This connects to the control as inference framework described in \cite{levine2018reinforcement} for the constrained RL formulation in \ref{eq:RL}. We consider a parameterized policy $\pi_{\theta}(y\mid x)$, aiming to minimize the divergence $\mathbb{D}_{\text{KL}}[\pi_{\theta}(y|x) \mid\mid \pi^*(y\mid x)]$, where $\pi^*$ denotes the optimal policy as characterized in Eq.~\ref{eq:optimum_model}, determined by the reward $r_{\phi}(y, x)$. Through straightforward manipulations, this leads to the following maximization objective:

\begin{equation}\label{eq:AC}
    \max_{\pi_{\theta}}\mathbb{E}_{\pi_{\theta}(y\mid x)}\bigg[\underbrace{r_{\phi}(x, y) -\beta\log\sum_{y}\pi_\text{ref}(y\mid x)\exp\left(\frac{1}{\beta}r_{\phi}(x, y)\right)}_{f(r_{\phi}, \pi_\text{ref}, \beta)} - \underbrace{\beta\log\frac{\pi_{\theta}(y\mid x)}{\pi_\text{ref}(y\mid x)}}_{\text{KL}}\bigg]
\end{equation}

The objective considered here is identical to that targeted by previous studies 
\citep{ziegler2020finetuning, stiennon2022learning, bai2022training, ouyang2022training}, where the DPO-equivalent reward is adopted for the reward class $r_{\phi}$. Within this framework, the normalization component in $f(r_{\phi}, \pi_\text{ref}, \beta)$ can be viewed as representing the soft value function for the reference policy $\pi_\text{ref}$. Although this normalization does not alter the location of the optimal solution, its exclusion may lead to high-variance policy gradients, thereby reducing the stability of the learning process. Employing a learned value function to estimate the normalization is possible, yet it often proves challenging to optimize. Alternatively, previous research has normalized the rewards relative to a human completion baseline, which effectively acts as a one-sample Monte Carlo approximation to the normalization factor. By contrast, the DPO reparameterization results in a reward structure that obviates the need for any such baselines.

\section{Experiments}
This section reports on the empirical assessment of DPO’s capacity to learn policies solely from preference data. We begin with controlled text-generation experiments to examine DPO's efficiency in balancing reward maximization with minimization of KL-divergence from the reference policy, in comparison to established preference optimization methods such as PPO. Subsequently, we analyze DPO's scalability and effectiveness on larger models and more challenging RLHF benchmarks, such as summarization and dialogue generation. Remarkably, DPO frequently matches or surpasses robust baselines, including RLHF with PPO and the best-of-$N$ sampling under a learned reward, often with minimal hyperparameter adjustment. Prior to discussing empirical outcomes, we outline the experimental methodology; further specifics are provided in Appendix~\ref{app:exp_details}.

\textbf{Tasks.} We evaluate our methods across three distinct open-ended text generation challenges. For each experiment, the algorithms are tasked with learning a policy from a preference dataset denoted by $\mathcal{D}=\bigl\{x^{(i)}, y_w^{(i)}, y_l^{(i)}\bigr\}_{i=1}^N$. In the case of \textbf{controlled sentiment generation}, $x$ represents the initial segment of a movie review sourced from the IMDb dataset \cite{maas-EtAl:2011:ACL-HLT2011}, with the objective for the policy to generate a continuation $y$ that exhibits positive sentiment. To facilitate systematic assessment, we synthetically create preference pairs for the generations using an existing sentiment classifier, ensuring that $p(\text{positive}\mid x,y_w)>p(\text{positive}\mid x,y_l)$. The SFT baseline is established by fine-tuning GPT-2-large on the training portion of IMDB reviews until optimal performance is attained (see further details in App~\ref{app:sentiment_details}). In the \textbf{summarization} task, $x$ corresponds to a Reddit forum post, while the policy must generate $y$, a summary capturing the post’s main ideas. This experiment employs the Reddit TL;DR dataset \citep{volske-etal-2017-tl}, complemented with human preference data collected by \citeauthor{stiennon2022learning}. For SFT, we utilize a model previously fine-tuned on human-generated summaries of forum posts\footnote{\url{https://huggingface.co/CarperAI/openai_summarize_tldr_sft}}, incorporating the TRLX \citep{leandro_von_werra_2023_7790115} RLHF library. Note that the human preference dataset from \citeauthor{stiennon2022learning} involves samples derived from another similarly-trained SFT model. The \textbf{single-turn dialogue} setting considers $x$ as a single human prompt—ranging from scientific queries to personal advice—requiring the policy to return a relevant and constructive reply $y$. For this, we rely on the Anthropic Helpful and Harmless dialogue dataset \citep{bai2022training}, which consists of 170,000 conversations between people and an AI assistant; each includes a pair of assistant replies and a human-annotated preference for the superior one. As there is no pre-existing SFT model in this scenario, we create one by fine-tuning a standard language model exclusively on the human-preferred completions.

\textbf{Evaluation.} We utilize two primary evaluation strategies in our experiments. For the controlled sentiment generation scenario, where the true reward function (represented by a sentiment classifier) is available, we assess each algorithm based on the trade-off between the maximized reward and the KL-divergence from the reference policy. Here, it is feasible to map out the reward–divergence frontier since the underlying reward is accessible. Conversely, in practical, real-world conditions where the true reward function is unknown, we employ the \textit{win rate} metric, comparing each algorithm to a baseline policy. In both the summarization and single-turn dialogue contexts, GPT-4 serves as a stand-in for human evaluators, assessing either the overall quality of summaries or the helpfulness of dialogue responses. For the summarization task, the baseline is drawn from reference summaries in the test dataset; for dialogue, the preferred response from the test set serves as the comparator. While prior research indicates that LMs may serve as more effective automated evaluators than existing quantitative measures \citep{Chen2023ExploringTU}, we further validate our reliance on GPT-4 in Section~\ref{sec:human-judgments} via a human evaluation study. Our findings show a high degree of alignment between GPT-4 and human raters, with agreement between humans and GPT-4 being at least as strong as, or exceeding, that observed between different human annotators.

\textbf{Methods.} Alongside DPO, we include several established methods for aligning language models with human preferences in our evaluation. As baseline prompting techniques, we apply zero-shot prompts using \textbf{GPT-J} \citep{gpt-j} in summarization and 2-shot prompts with \textbf{Pythia-2.8B} \citep{biderman2023pythia} in dialogue. We further assess the \textbf{SFT} model and introduce \textbf{Preferred-FT}, which is produced by supervised fine-tuning on the preferred completions $y_w$; these are sourced either from the SFT outputs (in sentiment and summarization) or from a generic language model (in the single-turn dialogue scenario). Another method in this category is \textbf{Unlikelihood}~\citep{welleck2019neural}, which adjusts the policy to simultaneously increase the likelihood of $y_w$ and decrease that of $y_l$, incorporating an optional scaling factor $\alpha \in [0,1]$ for the unlikelihood loss. We also evaluate \textbf{PPO} \citep{schulman2017proximal}, trained using a reward model derived from preference data, as well as \textbf{PPO-GT}, which is an oracle variant utilizing the ground-truth reward available in controlled sentiment experiments. For PPO-GT, we experiment with both an off-the-shelf version \cite{leandro_von_werra_2023_7790115} and an enhanced version featuring reward normalization and fine-tuned hyperparameters for improved outcomes; these improvements are also employed in the standard PPO experiments with learned reward models. Lastly, the \textbf{Best of $N$} approach is included as a strong baseline, where $N$ responses are generated from either the SFT model (or Preferred-FT in the case of dialogue), and the top response according to a preference-trained reward model is selected. Despite its efficacy, this method is computationally demanding—requiring $N$ samples for every test input—making it less feasible for practical deployment when $N$ is not small.

\subsection{How well can DPO optimize the RLHF objective?}

In standard RLHF algorithms, the reward maximization objective is regularized by a KL constraint to strike a balance between maximizing rewards and limiting policy drift from the reference model. Consequently, when evaluating and comparing different algorithms, it is critical to assess both the reward achieved and the associated KL divergence; achieving marginally higher rewards at the expense of substantially greater KL divergence is often suboptimal. Figure~\ref{fig:frontier-tldr-main} illustrates the trade-off frontier between reward and KL for a range of algorithms within the sentiment task. For each algorithm, we conduct several training runs, each with a distinct setting for policy conservativeness (target KL $\in\{3,6,9,12\}$ for PPO, $\beta \in \{0.05,0.1,1,5\}$, $\alpha\in\{0.05,0.1,0.5,1\}$ for unlikelihood, and different random seeds for preferred-FT). This grid search comprises 22 experimental runs in total. Every 100 training steps until the algorithms converge, policies are evaluated on a suite of held-out test prompts, reporting the average reward—computed via the true reward function—and the mean sequence-level KL\footnote{Here, sequence-level KL refers to the cumulative KL divergence across all timesteps.} with respect to the reference policy, $\text{KL}\left(\pi\mid \mid \pi_\text{ref}\right)$. The results demonstrate that DPO generates a highly efficient trade-off frontier, delivering superior rewards at consistently lower KL values. This finding stands out for several reasons. Notably, although both DPO and PPO are optimizing for the same underlying objective, DPO demonstrates significantly greater efficiency, outperforming PPO in terms of reward/KL trade-off. Additionally, DPO’s frontier is superior even compared to PPO when the latter has access to ground truth rewards (PPO-GT).

\subsection{Can DPO scale to real preference datasets?}
\label{sec:dpo-real-datasets}

We further investigate the fine-tuning capabilities of DPO in the domains of summarization and single-turn dialogue. In the case of summarization, widely-used automatic metrics such as ROUGE have been observed to correlate poorly with human judgment~\citep{stiennon2022learning}. Previous research has indicated that leveraging PPO for preference-based fine-tuning of language models leads to the generation of higher-quality summaries. To assess the comparative performance of different approaches, we generate completions using the test partition of the TL;DR summarization dataset and calculate their average win rates when matched against the reference completions in the test set. Sampling is performed at various temperatures ranging from 0.0 to 1.0, with the corresponding win rates presented in Figure~\ref{fig:frontier-tldr-main} (right). DPO, PPO, and Preferred-FT all build upon the same base GPT-J SFT model\footnote{\url{https://huggingface.co/CarperAI/openai_summarize_tldr_sft}}. Our experiments reveal that DPO achieves an approximate win rate of 61\% at temperature 0.0, outperforming PPO, which reaches around 57\% at its optimal temperature setting. Furthermore, the peak win rate obtained by DPO exceeds that of the best $N$ baseline. It is worth mentioning that we made no substantial efforts to optimize DPO's $\beta$ hyperparameter, suggesting that the present findings may not fully capture DPO’s optimal performance. Additionally, DPO demonstrates greater resilience to changes in sampling temperature compared to PPO, whose win rate can deteriorate to the level of the base GPT-J model as temperature increases. The Preferred-FT method fails to yield a meaningful improvement over the initial SFT model. For a direct human comparison, we further evaluate DPO and PPO in Section~\ref{sec:human-judgments}; in these assessments, human raters selected DPO samples generated at temperature 0.25 over PPO outputs produced at the same temperature 58\% of the time.

For single-turn dialogue experiments, we conduct assessments across various approaches using a segment of the test partition from the Anthropic HH dataset \citep{bai2022training}, specifically focusing on instances comprising a single human-assistant exchange. GPT-4 evaluations are based on the preferred completions from the test set, serving as references to determine the win rates for each method. In the absence of a widely-adopted SFT baseline for this task, we initialize with a pre-trained Pythia-2.8B model, utilize Preferred-FT to fine-tune a reference model on selected completions—ensuring outputs remain within the model's output distribution—and proceed to train with DPO. For comparison, we include the strongest result among 128 Preferred-FT completions (our analysis indicates that the Best of $N$ baseline reaches a performance plateau at $N=128$; further details in Appendix Figure~\ref{fig:best-of-n}) and also evaluate a 2-shot prompted variant of the Pythia-2.8B base model. Across these evaluations, DPO consistently matches or outperforms other methods when using the optimal temperature for each. We further analyze an RLHF model, trained with PPO on the Anthropic HH dataset\footnote{\url{https://huggingface.co/reciprocate/ppo_hh_pythia-6B}} using established implementations\footnote{\url{https://github.com/CarperAI/trlx/tree/main/examples/hh}}, but are unable to identify any prompting or sampling configurations that surpass the unmodified Pythia-2.8B. Given our TL;DR results and the alignment of reward objectives, we consider the Best of 128 benchmark a rough stand-in for PPO-level outcomes. In summary, DPO stands out as the sole efficient strategy that consistently enhances upon the preferred completions in the Anthropic HH dataset, while matching or outperforming the more resource-intensive Best of 128 approach. Figure~\ref{fig:dialogue-main} further illustrates that DPO rapidly attains its peak performance.

\subsection{Generalization to a new input distribution}

To further assess the robustness of PPO and DPO under distributional changes, we analyze their policies—originally trained on the Reddit TL;DR summarization task—by testing them on a distinct dataset: the test split of CNN/DailyMail news articles \citep{nallapati-etal-2016-abstractive}. Both methods are evaluated using their optimal sampling temperatures determined from TL;DR (0 and 0.25). The corresponding outcomes are shown in Table~\ref{tab:ood}. To measure comparative performance, we calculate GPT-4 win rates versus ground-truth summaries, employing the identical GPT-4 (C) prompt from the Reddit TL;DR experiments but substituting “forum post” with “news article.” On this unfamiliar distribution, DPO continues to achieve superior results relative to PPO. These findings offer preliminary evidence that DPO’s generalization capabilities can match or exceed those of PPO—even though DPO does not exploit the extra unlabeled Reddit TL;DR prompts incorporated by PPO.

\subsection{Validating GPT-4 judgments with human judgments}
\label{sec:human-judgments}

We carry out a human evaluation to test the consistency of GPT-4's judgments, using output from the TL;DR summarization experiments and employing two distinct GPT-4 prompts. The \textbf{GPT-4 (S)} (simple) prompt asks only which summary better covers the critical content in the post, whereas the \textbf{GPT-4 (C)} (concise) prompt additionally requests consideration of conciseness—prompted by the observation that, with the (S) prompt, GPT-4 often prefers longer, more repetitive outputs compared to human annotators. The full text of each prompt can be found in Appendix~\ref{app:prompts}. To ensure a range of sample qualities, we compare three systems: the highest-performing (DPO, temp. 0.25), the lowest-performing (PPO, temp. 1.0), and a mid-range model (SFT, temp. 0.25), all evaluated against the best-performing greedy PPO baseline. In both prompt scenarios, the rate at which GPT-4’s rankings coincide with human choices closely matches inter-human agreement rates, implying GPT-4 serves as an effective surrogate for human evaluation (for DPO and PPO-1, multiple human annotations are obtained; for others, rater availability limits annotation depth). On the whole, win rates based on the \textbf{GPT-4 (C)} prompt tend to align more closely with human preferences, which motivates its use for primary analyses in Section~\ref{sec:dpo-real-datasets}. Further information about the human evaluation protocol, including the user interface and the full roster of annotators, can be found in Appendix~\ref{app:human-study}.

\section{Discussion}
Preference-based learning presents a robust and scalable approach for developing language models that are both capable and aligned with user intentions. In this work, we introduced DPO, a straightforward methodology for leveraging preferences to train language models, circumventing the need for reinforcement learning. Rather than reformulating the preference learning problem to fit a conventional RL paradigm just to utilize standard RL methods, DPO establishes an explicit correspondence between policies of language models and reward functions. This alignment facilitates the direct optimization of language models according to human preferences through a basic cross-entropy loss, all without engaging reinforcement learning or compromising generality. With minimal hyperparameter adjustment, DPO demonstrates performance comparable to, or even surpassing, established RLHF approaches, including those based on PPO; as a result, DPO significantly lowers the entry threshold for preference-based language model training.

\textbf{Limitations \& Future Work.} Our findings open up numerous avenues for further investigation. A key question concerns the extent to which DPO-trained policies generalize to novel data distributions relative to approaches relying on explicit reward signals. Preliminary evidence suggests comparable generalization to models trained with PPO, but a broader evaluation is warranted. Another question is whether self-labeling strategies utilizing DPO-trained policies can capitalize on unannotated prompts as effectively as other approaches. Moreover, understanding the manifestations of reward over-optimization within the direct preference optimization framework, such as the modest decline observed in Figure~\ref{fig:dialogue-main}-right, requires further exploration. While our experiments encompass models up to 6B parameters, extending DPO to current state-of-the-art models that are several orders of magnitude larger represents a promising future research path. On the evaluation side, our analysis indicates that GPT-4-calculated win rates are influenced by prompt phrasing; thus, optimal methodologies for eliciting accurate assessments from automated evaluators merit further exploration. Finally, DPO's applicability is not confined to human preference-driven language model training but can also extend to other generative modeling domains.